% AIPR 2026 one-page abstract submission (Springer Meteor: https://meteor.springer.com/aipr2026).
% Results-neutral draft for the final paper.
\documentclass[runningheads]{llncs}

\usepackage[T1]{fontenc}
\usepackage{microtype}
\usepackage[hidelinks]{hyperref}

\begin{document}

\title{Label-Efficient Dark-Vessel Detection in SAR: Does SAR-Domain
Pretraining Outperform Optical and ImageNet Transfer Across ViT and CNN?}
\titlerunning{Label-Efficient Dark-Vessel Detection in SAR}

\author{John Roth\inst{1,2} \and Kyle Wagner\inst{1,3} \and Liv d'Aliberti\inst{1}}
\authorrunning{J. Roth et al.}

\institute{Johns Hopkins University, Baltimore MD, USA\\
\email{\{jroth32,kwagne35,odalibe1\}@jh.edu}
\and
The MITRE Corporation, McLean VA, USA\\
\email{johnroth@mitre.org}
\and
Boeing -- Phantom Works, St. Louis MO, USA\\
\email{kyle.w.wagner@boeing.com}}

\maketitle

\begin{abstract}
Detecting dark vessels---vessels lacking matched Automatic Identification
System (AIS) reports---in spaceborne synthetic aperture radar (SAR) plays an important role in supporting
maritime safety and enforcement, including monitoring activity potentially associated with sanctions evasion. 
Since human-verified examples are scarce, we study whether source-domain pretraining improves label efficiency on
xView3 (Sentinel-1 GRD, 10\,m pixel spacing) and whether the effect persists
across architectures. Our controlled comparison uses approximately
size-matched ViT-B/16 (86M parameters) and ConvNeXt-V2-Base (89M) tracks.
Within each track, we compare random initialization with
pretrained encoders from ImageNet, optical remote sensing, and SAR. The ViT
track uses AugReg, SatDINO (DINO on fMoW-RGB), and SARMAE (masked autoencoding
on SAR-1M); the ConvNeXt track uses FCMAE and BigEarthNet Sentinel-2 and Sentinel-1
models. Because the
checkpoints differ in pretraining recipe as well as source domain, domain
effects are interpreted within architecture. Matched-role cross-track
comparisons test whether the observed pattern persists across architectures.

All arms use the same [VH, VV, VH$-$VV] input, CenterNet-style, anchor-free Gaussian-heatmap point detector,
optimizer, schedule, scene-level splits, and seed. We fine-tune nested 10\%,
25\%, 50\%, and 100\% fractions of the training scenes in a stratified
150-scene study subset. We report seed-0 distance-matched F1 on held-out
scenes, plus dark-vessel recall and near-shore F1 on a disjoint, near-shore,
50-scene human-verified set used once for final evaluation. Because training
subsets are not restricted to near-shore scenes, this final measure tests
coastal generalization under distinct clutter conditions. All reported values are seed-0 point estimates.

\keywords{Synthetic aperture radar \and Dark-vessel detection \and Transfer
learning \and Label efficiency \and Vision transformers \and Convolutional
neural networks \and xView3-SAR.}
\end{abstract}

\end{document}
